% 第6章 6.7.1：割与指示向量。两个单位权三角形 {1,2,3}、{4,5,6} 由单位权边 3-4 相连，
% 顶点 7 以权 0.6 挂在顶点 2 上。(a) S={7}：cut=0.6（全部非空二分中最小）；
% (b) S={1,2,3,7}：cut=1。RatioCut：(a) 0.6*(1+1/6)=0.7，(b) 1*(1/4+1/3)=7/12≈0.583（最小）。
% 指示向量 s_i=1（i 属于 S，蓝）或 0（绿）；s^T L s 等于跨组边权之和（红色边）。
\begin{tikzpicture}[x=1mm,y=1mm,font=\figurefont\footnotesize,text=FigureInk,
  vS/.style={circle,draw=FigureStroke,fill=FigureBlueFill,line width=.8pt,minimum size=6mm,inner sep=0pt},
  vC/.style={circle,draw=FigureStroke,fill=FigureGreenFill,line width=.8pt,minimum size=6mm,inner sep=0pt},
  edge/.style={draw=FigureStroke,line width=.8pt},
  cutedge/.style={draw=FigureRed,line width=1.3pt},
  sval/.style={font=\FigureNoteFont\fontsize{7.5}{9}\selectfont,text=FigureMuted},
  note/.style={font=\FigureNoteFont\fontsize{7.5}{9}\selectfont,text=FigureMuted}]
\foreach \dx/\case in {0/0,56/1}{
 \begin{scope}[xshift=\dx mm]
  \ifnum\case=0
   \def\sa{vC}\def\sb{vC}\def\sc{vC}\def\sg{vS}
  \else
   \def\sa{vS}\def\sb{vS}\def\sc{vS}\def\sg{vS}
  \fi
  \node[\sa] (v1) at (10,4) {$1$};
  \node[\sb] (v2) at (10,19) {$2$};
  \node[\sc] (v3) at (21,11.5) {$3$};
  \node[vC] (v4) at (33,11.5) {$4$};
  \node[vC] (v5) at (44,19) {$5$};
  \node[vC] (v6) at (44,4) {$6$};
  \node[\sg] (v7) at (0,27) {$7$};
  \draw[edge] (v1)--(v2)--(v3)--(v1);
  \draw[edge] (v4)--(v5)--(v6)--(v4);
  \ifnum\case=0
   \draw[edge] (v3)--(v4);
   \draw[cutedge,line width=1pt] (v2)--(v7) node[midway,above right=-.5mm,note,text=FigureRed]{$0.6$};
   \node[anchor=west] at (-4,36) {(a) $S=\{7\}$，割权$0.6$};
  \else
   \draw[cutedge] (v3)--(v4) node[midway,above,note,text=FigureRed]{$1$};
   \draw[edge,line width=.55pt] (v2)--(v7) node[midway,above right=-.5mm,note]{$0.6$};
   \node[anchor=west] at (-4,36) {(b) $S=\{1,2,3,7\}$，割权$1$};
  \fi
 \end{scope}}
\end{tikzpicture}
